\section{Continuous geometry of certificates}
\label{sec:geometry}

The cluster problem in continuous global optimization concerns the boxes
that remain near an optimum. Its classical analyses give upper estimates
for specified subdivision methods
\citep{duKearfott1994Cluster,wechsungSchaberBarton2014Cluster,
kannanBarton2017ConstrainedCluster}.
Here we first bound the contribution of one valid test, so the resulting
lower bounds apply to every permitted cover and split tree. Near-optimal
covering counts also appear in certified Lipschitz optimization
\citep{bachocCesariGerchinovitz2021LipschitzCertificates}. The scale here
is determined by the relaxation gap. A quadratic gap gives a
$\sqrt{\epsilon+\eta}$ scale; it is not an assumption that the objective
is quadratic.

\subsection{Covering and arcsine lower bounds}

Let $N_\infty(Y,\delta)$ denote the minimum number of closed cubes of side
$\delta$ covering $Y$, with $N_\infty(\varnothing,\delta)=0$. A raw valid
cover under \cref{cert:eq-bound-gap}, or an owned certificate satisfying
\cref{cert:eq-owned-valid}, will be called an $\alpha$-valid family in
this subsection. Its size is denoted by $K$.

\begin{theorem}[Vertex covering bound]\label{geom:thm-cover-lower}
For every $\eta\geq0$ and every $\alpha$-valid family covering $E(\eta)$,
\begin{equation}\label{geom:eq-cover-lower}
 K\geq 2^{-n}N_\infty\!\left(E(\eta),
                         2\sqrt{(\epsilon+\eta)/\alpha}\right).
\end{equation}
No smoothness, convexity, or manifold hypothesis is needed.
\end{theorem}
\begin{proof}
If $y$ belongs to an owner with test $C$, validity gives
$(d_i^C(y))^2\leq q_C(y)\leq(\epsilon+\eta)/\alpha$ for each coordinate.
Choosing the nearer endpoint in every coordinate puts $y$ in a cube of
the stated side about a vertex of $C$. Its $2^n$ vertices cover the
family's contribution. For a raw cover use $C$ itself as its owner.
\end{proof}

For integration over a feasible manifold, the orientation and number of
coordinate sheets matter. Let $S$ be an embedded $C^1$ manifold of
dimension $1\leq d\leq n$. At $y\in S$, choose an $n\times d$ orthonormal
tangent matrix $U(y)$. For a $d$-coordinate subset $I$, put
$J_I(y)=|\det U_I(y)|$. Cauchy--Binet gives $\sum_I J_I^2=1$.
Assign $y$ to the lexicographically first maximizing $I$, giving Borel
sets $S^I$. Define
\[
 M^I(S;A)=\sup_{x\in\R^I}
       \#\{y\in S^I\cap A:y_I=x\},\qquad
 M(S)=\sum_{|I|=d}M^I(S;\R^n).
\]
The definition is independent of the orthonormal tangent basis. It counts
only the sheets on which that projection is chosen.

\begin{lemma}[Per-test arcsine estimate]\label{geom:lem-arcsine}
For every closed box $B$ and Borel set $A$,
\begin{equation}\label{geom:eq-arcsine}
 \int_{S\cap B\cap A}q_B^{-d/2}\,d\mathcal H^d
 \leq C_{n,d}\sum_{|I|=d}M^I(S;A\cap B),\qquad
 C_{n,d}=\left(\frac{\pi^2}{d}\right)^{d/2}
                                  \binom nd^{1/2}.
\end{equation}
The inequality is understood for nonnegative extended integrals.
\end{lemma}
\begin{proof}
On $S^I$, $J_I\geq\binom nd^{-1/2}$. Arithmetic--geometric mean gives
\[
 q_B^{-d/2}\leq d^{-d/2}\prod_{i\in I}
                  [(y_i-l_i)(u_i-y_i)]^{-1/2}.
\]
Apply the coordinate area formula on $S^I\cap B\cap A$. It changes the
integral with the factor $J_I$ into the integral over $y_I$, counted with
the number of points in each fibre. Bound that number by
$M^I(S;A\cap B)$. Each nondegenerate interval contributes
\[
 \int_{l_i}^{u_i}[(t-l_i)(u_i-t)]^{-1/2}\,dt=\pi.
\]
The factor $d^{-d/2}\pi^d\binom nd^{1/2}$ and summation over $I$ give
\cref{geom:eq-arcsine}. If a projected coordinate is fixed in $B$, the
portion of $S^I\cap B$ has zero $d$-measure: its Jacobian is bounded below
and its projected image has zero Lebesgue measure. Thus positive-measure
terms use only nondegenerate intervals, which also proves the claim for
degenerate boxes.
\end{proof}

\begin{theorem}[Stratum lower bound]\label{geom:thm-stratum}
Suppose $S\subset F$ and $M(S)<\infty$. Every $\alpha$-valid family
covering $S$ satisfies
\begin{equation}\label{geom:eq-stratum-lower}
 K\geq\frac{\alpha^{d/2}}{C_{n,d}M(S)}
              \int_S(m+\epsilon)^{-d/2}\,d\mathcal H^d.
\end{equation}
For $d=n$, this becomes
$K\geq(\alpha n/\pi^2)^{n/2}\int_S(m+\epsilon)^{-n/2}$.
These conclusions hold for owned subsets, so shared closed boundaries
of positive stratum measure cause no overcounting.
\end{theorem}
\begin{proof}
Validity gives $(m+\epsilon)^{-d/2}\leq
\alpha^{-d/2}q_C^{-d/2}$ on each owned contribution. Apply
\cref{geom:lem-arcsine} and sum. A cover can overlap, since summing
nonnegative contributions only increases their integral. For disjoint
owners the contributions partition $S$.
\end{proof}

Finite coordinate multiplicity follows locally from a $C^2$ graph, with
uniform constants for graphs of uniformly bounded second derivatives.
\Cref{app:geometry:lem-graphs} proves this fact. Some such control is
necessary: the planar union of $k$ parallel segments inside a fixed box
has zero curvature on every segment but an integral proportional to $k$.
An assertion using curvature alone would omit this number of sheets.

\subsection{A constrained covering comparison}

For the upper bounds, let $D$ be a cube of side $s_0$. Specify the original
constraint model by
\[
 F=\{z\in D\cap P_{\mathrm{ex}}:
          g_j(z)\leq0,\ h_k(z)=0\},\qquad
 v(z)=\max\{0,\max_j g_j(z),\max_k|h_k(z)|\}.
\]
Empty maxima are omitted. The closed set $P_{\mathrm{ex}}$ contains
exactly kept restrictions, and $R_B\subset B\cap P_{\mathrm{ex}}$.
The functions in $v$ are the original relaxed constraint functions.

\begin{assumption}[Second-order upper error and feasibility error bound]
\label{geom:ass-upper}
There are constants $\tau>0$, $\kappa\geq0$, $v_0>0$, and $L_f\geq0$
such that
\begin{align}
 f_B(z)&\geq f(z)-\tau w(B)^2,&
 v(z)&\leq\tau w(B)^2 &&(z\in R_B),\label{geom:eq-upper}\\
 \dist_2(z,F)&\leq\kappa v(z)
 &&&&(z\in D\cap P_{\mathrm{ex}},\ v(z)\leq v_0),\label{geom:eq-eb}
 \intertext{and $f$ is $L_f$-Lipschitz on $D$ in Euclidean distance.
Here $w(B)$ is the largest side length.}
 \Lambda&=\tau(1+L_f\kappa).\nonumber
\end{align}
\end{assumption}

The distance bound is for the full feasible set, including the root and
exact constraints. A constraint qualification asserted only at global
minimizers does not supply \cref{geom:eq-eb}. It must be established on
the region used by relaxed witnesses or assumed directly.

Let $T_{\mathrm{bis}}(\epsilon)$ be the number of processed nodes of
uniform $2^n$-ary dyadic subdivision, with $U=f_*$. Write
$s_j=s_0 2^{-j}$ and let $N_j(Y)$ count closed level-$j$ grid cubes
meeting $Y$. A $2^n$-way split can be realized by $n$ successive
coordinate splits; thus its leaf partition is an admissible binary tree
partition.

\begin{lemma}[Localization of unpruned boxes]\label{geom:lem-localize}
Under \cref{geom:ass-upper}, there are $j_0$ and $C_0$, independent of
$\epsilon$, such that
\begin{equation}\label{geom:eq-level-sum}
 T_{\mathrm{bis}}(\epsilon)\leq C_0+
   2^n5^n\sum_{\substack{j\geq j_0\\\Lambda s_j^2>\epsilon}}
                 N_j(E(\Lambda s_j^2-\epsilon)).
\end{equation}
In particular this tree terminates at every positive tolerance.
\end{lemma}
\begin{proof}
Choose $j_0$ so that $\tau s_j^2\leq v_0$ and
$\kappa\tau s_j^2\leq s_j$ for $j\geq j_0$. An unpruned cube $B$ has
$L(B)<f_*-\epsilon$, so some $z\in R_B$ satisfies
$f_B(z)<f_*-\epsilon$, even if the infimum is not attained. By the upper
error and the distance bound a nearest feasible $y$ obeys
\[
 |y-z|_2\leq\kappa\tau s_j^2\leq s_j,\qquad
 0\leq m(y)<\Lambda s_j^2-\epsilon.
\]
Hence $\Lambda s_j^2>\epsilon$. Assign $B$ to a grid cube containing
$y$. At most $5^n$ level-$j$ cubes lie within sup distance $s_j$ of a
point in that cube. Coarse levels give a fixed finite count, and each
unpruned cube generates $2^n$ children. This proves the displayed bound
and finiteness.
\end{proof}

\begin{theorem}[Running-supremum covering comparison]\label{geom:thm-profile}
Assume the bound gap \cref{cert:eq-bound-gap} and
\cref{geom:ass-upper}. Put
\begin{equation}\label{geom:eq-profile}
 \Phi_\alpha(\epsilon)=\sup_{\eta\geq0}
 N_\infty\!\left(E(\eta),2\sqrt{(\epsilon+\eta)/\alpha}\right).
\end{equation}
For sufficiently small $\epsilon>0$, with constants depending on the
fixed instance and scheme,
\begin{equation}\label{geom:eq-profile-comparison}
 2^{-n}\Phi_\alpha(\epsilon)
 \leq N_{\mathrm{cover}}\leq N_{\mathrm{rect}}
 \leq N_{\mathrm{tree}}\leq T_{\mathrm{bis}}
 \leq C_0+C(1+\log(1/\epsilon))\Phi_\alpha(\epsilon).
\end{equation}
If an owned $\alpha$-valid family covers $F$ and has size $K_F$, then
also $T_{\mathrm{bis}}\leq C_0+C(1+\log(1/\epsilon))K_F$.
\end{theorem}
\begin{proof}
\Cref{geom:thm-cover-lower} gives the lower bound at every $\eta$.
For a term in \cref{geom:eq-level-sum}, set
$\eta=\Lambda s_j^2-\epsilon$. A cube of side $s_j$ meets at most
$3^n$ closed grid cubes of that side. Refining a cube of side $a s_j$
to cubes of side $s_j$ takes at most $\max\{1,\lceil a\rceil\}^n$
cubes. Apply this with $a=2\sqrt{\Lambda/\alpha}$ to bound the level
count by a fixed multiple of $\Phi_\alpha(\epsilon)$. There are
$O(1+\log(1/\epsilon))$ nonempty levels in the sum.

For the last assertion, validity puts every point of
$E(\Lambda s_j^2-\epsilon)$ within
$\sqrt{\Lambda/\alpha}s_j$ of a test vertex. Each test has $2^n$
vertices; each such cube meets a bounded number of level-$j$ cells.
This bounds the same level sum by a fixed multiple of $K_F$.
It does not identify the owned family with a raw cover.
\end{proof}

Define $\psi_\alpha(\eta)=
N_\infty(E(\eta),2\sqrt{\eta/\alpha})$ for $\eta>0$. Then
\begin{equation}\label{geom:eq-running}
 2^{-n}\sup_{\eta\geq\epsilon}\psi_\alpha(\eta)
 \leq\Phi_\alpha(\epsilon)
 \leq\sup_{\eta\geq\epsilon}\psi_\alpha(\eta).
\end{equation}
Indeed the two scales differ by at most $\sqrt2$ when $\eta\geq\epsilon$,
and splitting each coordinate interval in two accounts for the factor
$2^n$. When $\eta<\epsilon$, set inclusion and monotonicity bound the
term by $\psi_\alpha(\epsilon)$.

The supremum cannot generally be replaced by its value at
$\eta=\epsilon$ with constants uniform over an instance class.
\Cref{app:geometry:ex-wells} gives a smooth one-dimensional family
with uniform relaxation and derivative bounds, a bounded single-scale
count, and certificate size at least a constant times the number of
near-optimal wells.

Under the pointwise gap, \cref{cert:eq-work,geom:thm-profile} show that
same-relaxation tightening can
improve on uniform subdivision by at most a logarithmic factor when
comparison is made to augmented event and reduction work. A statement
about nodes alone needs a bound on rounds per counted node; a statement
about evaluations uses \cref{cert:eq-eval}. The logarithmic comparison
does not assert that tightening takes a bounded number of rounds.

\subsection{When optimal-set dimension gives the rate}

Let $M=E(0)$. An optimal set does not determine the rate by itself:
flat objective growth can make an isolated minimizer expensive.
The next statement makes the needed growth condition explicit.

\begin{proposition}[Quadratic growth and covering dimension]
\label{geom:prop-dimension}
Under the hypotheses of \cref{geom:thm-profile}, assume
\begin{equation}\label{geom:eq-qg}
 m(y)\geq c_g\dist_\infty(y,M)^2\qquad(y\in F),\qquad c_g>0.
\end{equation}
Then
\begin{equation}\label{geom:eq-qg-bounds}
 2^{-n}N_\infty(M,2\sqrt{\epsilon/\alpha})
 \leq N_{\mathrm{cover}}(\epsilon),\qquad
 T_{\mathrm{bis}}(\epsilon)\leq C_0+
     C\sum_{\Lambda s_j^2>\epsilon}N_\infty(M,s_j).
\end{equation}
If $M$ has box-counting dimension $p>0$, each of the three certificate
minima and $T_{\mathrm{bis}}$ has logarithmic growth exponent $p/2$:
$\log N(\epsilon)/\log(1/\epsilon)\to p/2$.
If $M$ is finite, the upper bound is $O(1+\log(1/\epsilon))$.
The upper assertion requires only \cref{geom:ass-upper} and quadratic
growth; it does not require an objective lower gap.
\end{proposition}
\begin{proof}
The lower bound takes $\eta=0$ in \cref{geom:thm-cover-lower}.
By quadratic growth,
$E(\Lambda s_j^2)\subset\{\dist_\infty(\cdot,M)
\leq\sqrt{\Lambda/c_g}s_j\}$. A cover of $M$ by $s_j$-cubes
therefore bounds the level count in \cref{geom:eq-level-sum} by a fixed
multiple of $N_\infty(M,s_j)$. For every $\delta>0$, the definition of
box dimension bounds this count above by $s_j^{-p-\delta}$ and below
by $s_j^{-p+\delta}$ at sufficiently small scales. The geometric upper
sum stops at $s_j\asymp\sqrt\epsilon$. Taking logarithms and letting
$\delta\downarrow0$ proves the exponent statement. A finite set has a
uniformly bounded covering count at every scale.
\end{proof}

\begin{corollary}[An optimal manifold]\label{geom:prop-manifold}
Under the hypotheses of \cref{geom:prop-dimension}, suppose $M$ is a
compact embedded $C^1$ manifold of dimension $p\geq1$, possibly with
boundary, and with a finite regular coordinate-chart cover. Then
\[
 N_{\mathrm{cover}}\asymp N_{\mathrm{rect}}
 \asymp N_{\mathrm{tree}}\asymp T_{\mathrm{bis}}
 \asymp\epsilon^{-p/2}.
\]
\end{corollary}
\begin{proof}
Finite bounded-slope charts give $N_\infty(M,s)\leq Cs^{-p}$.
A single relatively open chart contains a $p$-dimensional coordinate
cube; a sup-norm cube of side $s$ projects to a cube of that side, so
$N_\infty(M,s)\geq cs^{-p}$. Apply \cref{geom:eq-qg-bounds}; the
geometric sum is $O(\epsilon^{-p/2})$ for $p>0$.
\end{proof}

A stratum integral can give the complete rate under direct geometric
regularity conditions. Suppose finitely many strata $S_k\subset F$,
of dimensions $d_k\geq1$, have finite $M(S_k)$. For fixed positive
$t_0,r_0,\theta$ and finite constants $K_R,K$, assume that for
$0<t\leq t_0$: every $y\in E(t)$ is within
$\sqrt{K_Rt}$ in sup norm of a point $z\in S_k$ with $m(z)\leq K_Rt$;
on each stratum,
\[
 m(x)\leq K(m(z)+|x-z|_\infty^2),\qquad
 \mathcal H^{d_k}(S_k\cap\{ |x-z|_\infty\leq r\})
       \geq\theta r^{d_k}
\]
whenever $m(z)\leq K_Rt_0$, $|x-z|_\infty\leq r_0$, and
$0<r\leq r_0$. These are proximity, quadratic doubling, and lower
density hypotheses, respectively.

\begin{theorem}[Constrained stratum integral comparison]
\label{geom:thm-integral}
Under these hypotheses and those of \cref{geom:thm-profile}, put
$I_k(\epsilon)=\int_{S_k}(m+\epsilon)^{-d_k/2}\,d\mathcal H^{d_k}$.
Then
\[
 1+\sum_k I_k(\epsilon)\asymp N_{\mathrm{cover}}
 \asymp N_{\mathrm{rect}}\asymp N_{\mathrm{tree}}
 \asymp T_{\mathrm{bis}}.
\]
If some strata are finite sets, their contribution to the upper bound
is instead $O(\#S_k(1+\log(1/\epsilon)))$; the statement does not supply
a matching logarithmic lower bound from those points alone.
\end{theorem}
\begin{proof}
\Cref{geom:thm-stratum} bounds the count below by each $I_k$ and hence
by their sum up to a fixed factor. For the upper bound, at level $j$
choose a maximal sup-norm $s_j$-separated family in
$\{z\in S_k:m(z)\leq K_R\Lambda s_j^2\}$. Proximity and
\cref{geom:lem-localize} assign every unpruned cube to one of these
points within $Cs_j$, and each point receives at most a fixed number
of grid cubes. Their relative balls of radius $s_j/3$ are disjoint.
Density gives mass at least $\theta(s_j/3)^{d_k}$ in each; doubling puts
all this mass in $\{m+\epsilon\leq C's_j^2\}$. Thus the level count
is at most
$C\sum_k s_j^{-d_k}\mathcal H^{d_k}
(S_k\cap\{m+\epsilon\leq C's_j^2\})$.
Tonelli's theorem changes the sum over levels to an integral of
$\sum_{s_j\geq\sqrt{(m+\epsilon)/C'}}s_j^{-d_k}$.
For $d_k\geq1$ this geometric sum is bounded by
$C(m+\epsilon)^{-d_k/2}$. Coarse levels add a constant. A finite
zero-dimensional stratum gives at most its cardinality at each level,
which explains its stated contribution.
\end{proof}

\subsection{Gaps caused by relaxed constraints}

An objective can be kept exactly while relaxed constraints still force
large certificates. Feasible-point objective gaps do not describe this
case. Its witnesses can be infeasible points with objective below $f_*$,
the mechanism already identified in constrained cluster analysis
\citep{kannanBarton2017ConstrainedCluster}. For a lower bound we need
such points to be admitted by the specified relaxation.

\begin{assumption}[Inner tube]\label{geom:ass-tube}
For fixed $\alpha\geq0$ and $\beta>0$, every test box satisfies
\begin{equation}\label{geom:eq-tube}
 z\in B\cap P_{\mathrm{ex}},\quad v(z)\leq\beta q_B(z)
 \quad\Longrightarrow\quad
 z\in R_B,\quad f_B(z)\leq f(z)-\alpha q_B(z).
\end{equation}
\end{assumption}

Exactly enforced constraints must be included in $P_{\mathrm{ex}}$,
since \cref{geom:eq-tube} cannot promise admission of their violators.
The case $\alpha=0$ permits an exact objective.

\begin{lemma}[Owned tube dichotomy]\label{geom:lem-tube-events}
Every raw valid certificate under \cref{geom:ass-tube} satisfies
\begin{equation}\label{geom:eq-dichotomy}
 v(z)>\beta q_C(z)\quad\hbox{or}\quad
 f(z)-\alpha q_C(z)\geq f_*-\epsilon
                  \qquad(z\in C\cap P_{\mathrm{ex}}).
\end{equation}
For runs using splits, same-relaxation reductions, and raw or expanded
probe pruning, the ownership construction in \cref{cert:lem-events}
gives the same dichotomy for $z\in A\cap P_{\mathrm{ex}}$.
All terminal owners, including owners with no feasible point, have count
\begin{equation}\label{geom:eq-all-owners}
 K_{\mathrm{all}}\leq\ell_{\mathrm{aug}}+2nR_{\mathrm{rel}}.
\end{equation}
Original-feasibility reductions and other independent ways to remove
infeasible tube witnesses are excluded in this run statement.
\end{lemma}
\begin{proof}
If $v(z)\leq\beta q_C(z)$, then for a containing source $B$,
$q_C(z)\leq q_B(z)$ puts $z$ in its inner tube. Hence $z\in R_B$ and
$f_B(z)\leq f(z)-\alpha q_B(z)\leq f(z)-\alpha q_C(z)$.
A successful source bound gives the second alternative. An infeasible
source cannot contain such a tube point. For a discarded
same-relaxation owner, the removal predicate gives
$f_B(z)>U-\epsilon\geq f_*-\epsilon$ instead. Emptying reductions use
the same argument on their parent. The owners partition all of $D$,
and their count is the one in \cref{cert:lem-events} before filtering by
feasibility. Filtering is not allowed here because the witnesses can
be infeasible.
\end{proof}

For $\nu>0$ define
$W(\epsilon,\nu)=\{z\in D\cap P_{\mathrm{ex}}:
v(z)\leq\nu,\ f(z)<f_*-\epsilon\}$.
Every family of $K$ test-owner pairs covering these points and satisfying
\cref{geom:eq-dichotomy} obeys
\begin{equation}\label{geom:eq-infeasible-cover}
 K\geq2^{-n}N_\infty(W(\epsilon,\nu),2\sqrt{\nu/\beta}).
\end{equation}
Indeed the objective alternative is impossible, so each coordinate's
distance to the nearer test endpoint is less than $\sqrt{\nu/\beta}$.

\begin{theorem}[Compatible nonlinear constraint transfer]
\label{geom:thm-tube-kkt}
Let $z_*$ be a global minimizer. On an open neighborhood of $z_*$,
suppose $f$, a finite constraint description of the unchanged exact set
$P_{\mathrm{ex}}$, and every original relaxed function defining $v$
are $C^2$. Combine that description, the root inequalities, and these
original relaxed functions. Assume LICQ for all equalities and active
inequalities in this combined tuple and KKT multipliers for the same
tuple. Suppose an active relaxed inequality has a positive multiplier,
or a relaxed equality has a nonzero multiplier.

There are a neighborhood $U_*$ and constants $c,\epsilon_0>0$ such that
for $Y\subset E(\eta)\cap U_*$ and $0<\eta+\epsilon\leq\epsilon_0$,
every dichotomy family covering $D\cap P_{\mathrm{ex}}$ satisfies
\begin{equation}\label{geom:eq-local-tube-cover}
 K\geq2^{-n}N_\infty(Y,c\sqrt{\eta+\epsilon}).
\end{equation}
If the manifold on which all active constraints equal their values at
$z_*$ has dimension $d\geq1$, then on a fixed relatively open feasible
patch $S$ of that manifold,
\begin{equation}\label{geom:eq-tube-integral}
 K\geq c_1\int_S(m+\epsilon)^{-d/2}\,d\mathcal H^d
             \geq c_2\log(1/\epsilon)
\end{equation}
for sufficiently small $\epsilon>0$. Neither strict complementarity
of all active inequalities nor a second-order sufficient condition is
needed for this lower bound.
\end{theorem}

The complete proof in \cref{app:geometry:proof-tube-kkt} constructs
constraint-preserving curves by the implicit function theorem. It proves
uniform injectivity, coordinate multiplicity, and area bounds for their
tolerance-dependent shifted patches before applying the arcsine estimate.
The compatibility requirement preserves the original exact/relaxed
assignment and violation function. A LICQ description of $F$ alone
cannot replace it. For example, let $P_{\mathrm{ex}}$ be the unit ball,
relax $1-|z|^2\leq0$, and minimize $-z_1$ over their intersection.
The feasible set is the unit sphere and its equality description is
LICQ at $e_1$, but the original exact and relaxed gradients are opposite.
With $R_B=B\cap P_{\mathrm{ex}}$ and the exact objective, the root is
already exact and satisfies the tube condition with $\alpha=0$.
No point in $P_{\mathrm{ex}}$ has objective below $-1$. The combined
LICQ hypothesis correctly excludes this instance.

For a run, \cref{geom:eq-tube-integral} applies to $K_{\mathrm{all}}$,
and hence to augmented event and reduction work by
\cref{geom:eq-all-owners}. Under the distinct-evaluation charging model,
retain every successful source, including infeasibility sources, and
the discarded owners. They cover $D\cap P_{\mathrm{ex}}$ and satisfy
the dichotomy, giving a fresh count
$K_{\mathrm{eval,all}}\leq(2n+1)S_{\mathrm{eval}}$.
Thus the same lower bound applies to evaluations in that model.
Exact propagation of original constraints can remove the infeasible
witnesses and is outside these conclusions.

\begin{corollary}[Finite regular KKT minimizers]\label{geom:cor-kkt-rates}
Assume \cref{geom:ass-upper} and that the set of global minimizers is
finite. Near each minimizer, suppose $f$ and a finite description of the
full feasible set are $C^2$, the description satisfies LICQ and KKT,
every active inequality has
a positive multiplier, and the Lagrangian Hessian is positive definite
on the common nullspace of the active inequality and equality gradients.
Then $T_{\mathrm{bis}}=O(1+\log(1/\epsilon))$.

If at least one active stratum has dimension at least one and the
objective bound gap holds, all three raw certificate minima and
$T_{\mathrm{bis}}$ are $\Theta(\log(1/\epsilon))$. The same matching
rate holds under the tube hypothesis when that minimizer satisfies
the compatible original-tuple and relaxed-multiplier hypotheses of
\cref{geom:thm-tube-kkt}.

If instead there is one minimizer $z_*$ and its active stratum has
dimension zero, a finite exact split-tree certificate exists under the
stronger vertex-vanishing upper error for some $\alpha'\geq0$,
\begin{equation}\label{geom:eq-upper-q}
 f_B(z)\geq f(z)-\alpha' q_B(z),\qquad
                    v(z)\leq\alpha'q_B(z)\qquad(z\in R_B).
\end{equation}
Thus $N_{\mathrm{tree}}(\epsilon)=O(1)$ even at $\epsilon=0$.
\end{corollary}

\Cref{app:geometry:proof-kkt-rates} proves quadratic growth from these
KKT assumptions, the local stratum logarithm, and sharp growth at a
zero-dimensional stratum. The error bound for the relaxation remains
an independent assumption. For the feasible-objective branch, a local
description of the full feasible set suffices for these geometric facts;
the tube branch specifically requires compatibility with the original
exact/relaxed tuple.

Width-squared upper convergence alone does not give the last conclusion.
For $F=[0,1]$, $f(x)=x$, and $f_B(x)=x-\tau w(B)^2$ on $R_B=B$, the
minimum at zero is sharp, but every positive-width interval containing
zero has a negative bound. A finite closed interval cover of $[0,1]$
must contain such an interval; singleton tests at zero cannot cover all
small positive points. Hence it has no finite exact certificate.
Under a positive objective lower gap, uniform bisection still requires
$\Omega(\log(1/\epsilon))$ when a minimizer coordinate stays in a fixed
middle portion of its dyadic interval at every level. A root-coordinate
position with an odd rational denominator greater than one has this
property. Merely being nondyadic does not imply it.

\begin{example}[Isotropic and directional constraint gaps]
\label{geom:ex-flat-tube}
Let $D=[-6/5,13/10]^2$, $f(z)=z_2$, and let $F$ be given by
$-z_2\leq0$, with $P_{\mathrm{ex}}=\R^2$. Keep $f$ exact.
For $R_B=\{z\in B:-z_2\leq q_B(z)\}$, the tube coefficient is $\beta=1$.
The strip $[-6/5,13/10]\times[-2\epsilon,-\epsilon)$ lies in
$W(\epsilon,2\epsilon)$ for small $\epsilon$. Its horizontal projection
has length $5/2$, so \cref{geom:eq-infeasible-cover} gives
$N_{\mathrm{cover}}\geq c\epsilon^{-1/2}$. The upper error,
$\dist(z,F)=v(z)$, and quadratic growth to the optimal segment give the
upper part of \cref{geom:prop-dimension} without an objective lower gap.
The segment has $N_\infty(M,s)=O(s^{-1})$, so the geometric level sum
is $O(\epsilon^{-1/2})$. Combining the upper and infeasible-cover lower
bounds gives all three raw minima $\Theta(\epsilon^{-1/2})$.

For the directional relaxation
$R_B=\{z\in B:-z_2\leq a_2^B(z_2)\}$, three horizontal slabs give an
exact certificate: $[-6/5,-3/5]$, $[-3/5,0]$, and $[0,13/10]$ in
the second coordinate. The first is infeasible since
$-z_2\geq3/5>9/100\geq a_2^B$; in the middle slab a negative feasible
$z_2$ would require $1\leq z_2+3/5$, which is impossible; the last has
minimum zero. This scheme has no isotropic tube coefficient $\beta>0$.
Finally, exact propagation of the original inequality retains $z_2\geq0$
and makes even the isotropic relaxation exact in one node.
\end{example}

\subsection{McCormick gaps and graph geometry}

The objective-gap hypothesis \cref{cert:eq-bound-gap} requires loss in
every coordinate not fixed by a box vertex. Bilinear envelopes have a
different zero set. In this subsection, fix the decomposition
\begin{equation}\label{geom:eq-mcc-model}
 f(x)=g(x)+\sum_{ij\in G}c_{ij}x_ix_j,
 \qquad F=D\cap P,
\end{equation}
where $g$ is convex and continuous, $P$ is a polyhedron in original
variables, $G$ has an edge $ij$ for every nonzero coefficient, and each
bilinear term is relaxed separately by its convex envelope on $B$.
The polyhedron is kept exactly. Write
$f_B=f-\Gamma_B$ for this relaxation, so
\begin{equation}\label{geom:eq-mcc-valid}
 L(B)\geq f_*-\epsilon\quad\Longleftrightarrow\quad
 \Gamma_B(x)\leq m(x)+\epsilon\quad(x\in F\cap B).
\end{equation}
The fixed decomposition matters even when the full objective is convex.
For an owned test $(C,A)$ in this subsection, \emph{validity} means
$\Gamma_C(x)\leq m(x)+\epsilon$ for $x\in A\cap F$. This is the
McCormick validity condition, rather than the isotropic $q$ condition
used earlier. Every owned-family statement below requires this condition.
The formulas below are the classical bilinear-envelope formulas
\citep{mccormick1976Computability}; we give their derivation to identify
the geometric properties used by the certificate bounds.

\begin{lemma}[Bilinear gaps and exact faces]\label{geom:lem-mcc-gap}
Put $\delta_i^-=x_i-l_i$ and $\delta_i^+=u_i-x_i$ on $B$. The gap for
the term $c_{ij}x_ix_j$ is
\begin{equation}\label{geom:eq-edge-gap}
 \gamma_{ij,B}(x)=|c_{ij}|\begin{cases}
 \min(\delta_i^-\delta_j^-,\delta_i^+\delta_j^+),&c_{ij}>0,\\
 \min(\delta_i^-\delta_j^+,\delta_i^+\delta_j^-),&c_{ij}<0.
 \end{cases}
\end{equation}
It satisfies
\begin{align}
 |c_{ij}|d_i^B d_j^B&\leq\gamma_{ij,B}
       \leq |c_{ij}|\min(d_i^Bw_j,d_j^Bw_i),\label{geom:eq-edge-bounds}\\
 \gamma_{ij,B}&\leq\tfrac12|c_{ij}|(a_i^B+a_j^B),&
 \max_B\gamma_{ij,B}&=\tfrac14|c_{ij}|w_iw_j.\label{geom:eq-edge-upper}
\end{align}
Moreover $\Gamma_B=\sum_{ij\in G}\gamma_{ij,B}$ vanishes at $x$
precisely when the coordinates of $x$ at box endpoints form a vertex
cover of $G$. Thus it vanishes on an entire coordinate face precisely
when that face's free coordinates form an independent set.
\end{lemma}
\begin{proof}
Normalize the two nondegenerate intervals to $s,t\in[0,1]$, removing
affine terms. The convex envelope of $st$ is $\max(0,s+t-1)$.
For $s+t\leq1$, express $(s,t)$ as a convex combination of
$(0,0),(1,0),(0,1)$, whose weighted product is zero. For $s+t\geq1$,
use $(1,0),(0,1),(1,1)$, whose weighted product is $s+t-1$.
Every convex underestimator is bounded above by these combinations,
and the two affine minorants are valid. Similarly the concave envelope
is $\min(s,t)$, using, when $s\leq t$, weights $s,t-s,1-t$ on
$(1,1),(0,1),(0,0)$ and reversing coordinates otherwise. Scaling and
the sign of the coefficient give \cref{geom:eq-edge-gap}.

Each endpoint distance is at least its nearest-endpoint distance;
choosing the product involving the nearest endpoint in one coordinate
gives \cref{geom:eq-edge-bounds}. The product of the two candidate
products in the minimum is $a_i^Ba_j^B$, so the minimum is at most
$\sqrt{a_i^Ba_j^B}\leq(a_i^B+a_j^B)/2$. It is also at most
$w_iw_j/4$, with equality at the common interval centers.
Degenerate intervals give zero directly. Each edge gap is nonnegative
and vanishes exactly when one endpoint coordinate is fixed; their sum
therefore has the claimed graph zero set.
\end{proof}

The upper inequalities give
\begin{equation}\label{geom:eq-mcc-upper}
 \Gamma_B\leq\tfrac12\left(\max_i\sum_{j:ij\in G}|c_{ij}|\right)q_B,
 \qquad
 \max_B\Gamma_B=\tfrac14\sum_{ij\in G}|c_{ij}|w_iw_j.
\end{equation}
All edge maxima occur at the same box center. The first inequality
supplies a vertex-vanishing upper error, but supplies no positive
isotropic objective lower gap.

A joint envelope of the bilinear sum is pointwise stronger than the
sum of separate envelopes. It has the same exact-face criterion:
any convex representation of a point on a box face uses points of
that face, by a supporting affine functional. Its restriction is
therefore the envelope of the restricted sum. That restricted bilinear
sum is convex exactly when it is affine: a surviving free edge has
the indefinite Hessian submatrix $\left(\begin{smallmatrix}0&c\\c&0
\end{smallmatrix}\right)$. Sums of termwise edge losses below are not
claims about a joint envelope.

The ownership and cost construction also applies to this model.
Indeed, restricting a convex underestimator from $B$ to $C\subset B$
gives an admissible underestimator on $C$, so each child envelope is at
least the restricted parent envelope. Hence $\Gamma_C\leq\Gamma_B$
on $C$. At a discarded same-relaxation feasible point,
$f_B(x)>U-\epsilon$ gives $\Gamma_B(x)<m(x)+\epsilon$, and therefore
the owned child test is McCormick-valid. Successful containing source
bounds give the weak inequality in the same way. Original-feasibility
owners contain no feasible point. Thus the rectangular ownership proof
gives $K_F\leq\ell_{\mathrm{aug}}+2nR_{\mathrm{rel}}$ for this
validity condition as well. The augmented-node and separately charged
evaluation conclusions in \cref{cert:eq-work,cert:eq-eval} follow with
the same distinct costs. They do not require an isotropic objective gap.

\begin{theorem}[Flat near-optimal sets]\label{geom:thm-flat-graph}
Suppose a selected nonempty graph $H\subset G$ satisfies
$|c_{ij}|\geq\alpha>0$ on every edge. Let $V$ have orthonormal columns
spanning a $p$-flat, $p\geq1$, and let $v^i$ be its rows. For a
full-dimensional bounded convex body $K\subset\R^p$ put
$W_i(K)=\sup_{t\in K}v^i\cdot t-\inf_{t\in K}v^i\cdot t$ and define
\[
 \tau(V,H)=\inf_K
   \frac{(\max_{ij\in H}W_i(K)W_j(K))^{p/2}}{\vol_p(K)}.
\]
For every bounded convex $S\subset F\cap(z_0+\operatorname{range}V)$
with $m\leq\eta$ on $S$, every raw valid cover or rectangular-owner
family covering $S$ has size
\begin{equation}\label{geom:eq-flat-graph}
 K\geq\tau(V,H)
   \left(\frac{\alpha}{(p+1)^2(\eta+\epsilon)}\right)^{p/2}
                                                     \mathcal H^p(S).
\end{equation}
If every vertex-cover projection is injective on $\operatorname{range}V$,
then $\tau(V,H)>0$. This sufficient condition is called graph
transversality.
\end{theorem}
\begin{proof}
Intersect $S$ with an owner $A\subset C$. For a positive-volume
intersection its centroid belongs to that convex intersection even
when some endpoints are open. The elementary supporting-plane estimate
proved in \cref{app:geometry:lem-centroid} puts its nearest distance to
each pair of box faces at least $W_i/(p+1)$. At that owned centroid,
validity and \cref{geom:eq-edge-bounds} give
$\max_{ij\in H}W_iW_j\leq(p+1)^2(\eta+\epsilon)/\alpha$.
The definition of $\tau$ bounds the intersection volume; summing
over the family proves \cref{geom:eq-flat-graph}.

For transversality define
\[
 \Delta=\min_{J_0\text{ vertex cover of }H}
           \max_{\substack{J\subset J_0\\|J|=p}}|\det V_J|.
\]
Injectivity makes $\Delta>0$. For any body $K$, put
$b=\max_{ij\in H}W_iW_j$. The coordinates with $W_i\leq\sqrt b$
form a vertex cover. Among them choose $p$ rows with determinant at
least $\Delta$. Their coordinate ranges enclose $K$ in a parallelepiped
of volume at most $b^{p/2}/\Delta$. Thus $\tau(V,H)\geq\Delta$.
\end{proof}

For a unit line direction $v$, the formula reduces to
$\tau(v,H)=\max_{ij\in H}\sqrt{|v_iv_j|}$. In higher dimension
transversality is stronger than positivity of $\tau$: on $\R^2$ with
one edge, $\tau=1$ because area is at most the product of coordinate
widths, yet projecting to one endpoint is not injective on $\R^2$.

\begin{proposition}[Curved transversal strata]\label{geom:prop-curved-graph}
Let $S\subset F$ be a compact subset of an embedded $C^1$
$p$-manifold, and suppose every tangent space is graph-transversal for
the selected graph $H$ above. If $m\leq\eta$ on $S$, then for small
$\eta+\epsilon$ every valid family covering $S$ satisfies
\[
 K\geq c_S\mathcal H^p(S)
                           (\alpha/(\eta+\epsilon))^{p/2}.
\]
In particular an optimal stratum of positive measure forces
$\Omega(\epsilon^{-p/2})$ tests.
\end{proposition}
\begin{proof}
For each vertex cover $J_0$, injectivity on each tangent space gives
locally a $p$-coordinate projection with coordinates in $J_0$ and an
inverse chart of bounded derivative. A finite chart cover gives
\[
 \mathcal H^p(S\cap\{|x_{J_0}-\xi|_\infty\leq r\})\leq Cr^p
\]
uniformly in $\xi$ and small $r$. This follows by projecting each chart
intersection into the $p$ chosen coordinate intervals and bounding
the inverse chart's area factor.
At an owned point of a valid test, put
$\rho=\sqrt{(\eta+\epsilon)/\alpha}$. The edge bounds imply
$d_i d_j\leq\rho^2$ on $H$. Consequently $\{i:d_i\leq\rho\}$ is
a vertex cover and contains a minimal cover $J_0$. The point's
$J_0$ coordinates lie within $\rho$ of one of that test's
$2^{|J_0|}$ endpoint combinations. Sum the slab measure estimate over
the finitely many minimal covers and endpoint choices. Each test
contributes at most $C'\rho^p$ measure, proving the assertion.
\end{proof}

\begin{proposition}[Fractional vertex-cover gap budget]
\label{geom:prop-fractional}
Suppose an $n$-dimensional feasible cube $R$ of side $W$ has
$m\leq\eta$ throughout, and $H\subset G$ has coefficients at least
$\alpha>0$ in magnitude. Let
\[
 \tau^*(H)=\min\left\{\sum_i z_i:z_i\geq0,
                                  \ z_i+z_j\geq1\ (ij\in H)\right\}.
\]
If $\rho^2=4(\eta+\epsilon)/\alpha\leq W^2$, every raw valid cover
of $R$ has size at least
\begin{equation}\label{geom:eq-frac-lower}
 K\geq(W/\rho)^{2\tau^*(H)}.
\end{equation}
Conversely, on a root cube $D$ of side $W$, for the full relaxed graph
$G$ there is a grid tree with
\begin{equation}\label{geom:eq-frac-upper}
 K\leq2^n(W/\rho')^{2\tau^*(G)},\qquad
 0<\rho'\leq W,\quad
 \frac{\rho'^2}{4}\sum_{ij\in G}|c_{ij}|\leq\epsilon.
\end{equation}
This is a statement about the specified gap budget and near-optimal
cube. It does not determine the complexity from the optimal set alone.
\end{proposition}
\begin{proof}
The center of $C\cap R$, when this intersection has positive volume,
is at distance at least half each intersection width $w_i$ from the
corresponding faces of $C$. Validity gives $w_iw_j\leq\rho^2$ on $H$.
Thus its volume is at most
\[
 \max\left\{\prod_iw_i:0<w_i\leq W, w_iw_j\leq\rho^2\ (ij\in H)
                                                    \right\}.
\]
For $W>\rho$, substitute $w_i=W(\rho/W)^{2z_i}$. The constraints
become exactly those defining $\tau^*(H)$, and the maximum product
is $W^n(\rho/W)^{2\tau^*(H)}$. The equality case $W=\rho$ is direct.
Summing intersection volumes proves \cref{geom:eq-frac-lower}.

Choose an optimal fractional cover $z$ of $G$ with
$z_i\in\{0,1/2,1\}$; \cref{app:geometry:lem-half-integral} proves
its existence. Grid side lengths
$b_i=W(\rho'/W)^{2z_i}$ satisfy $b_ib_j\leq\rho'^2$ on every
relaxed edge. \Cref{geom:eq-mcc-upper} bounds the gap in every grid
cell by $\epsilon$, so $m\geq0$ makes every cell valid. The grid is
realizable by coordinate cuts, and its size is at most
$\prod_i\lceil W/b_i\rceil\leq2^n(W/\rho')^{2\tau^*(G)}$.
The upper construction must use every relaxed edge; a subgraph used
for the lower bound cannot omit edges from its gap budget.
\end{proof}

\begin{proposition}[Concave matching slices]\label{geom:prop-matching}
Let $\{i_rj_r:r=1,\ldots,k\}$ be a matching in $G$, with $k\geq1$, and put
$v_r=e_{i_r}-\operatorname{sign}(c_{i_rj_r})e_{j_r}$,
$Vt=\sum_r t_rv_r$, and $T_z=\{t:z+Vt\in F\}$.
Every valid raw cover, or owned valid family covering this slice, obeys
\begin{equation}\label{geom:eq-matching}
 K\geq\left(\frac{k}{\pi^2}\right)^{k/2}
       \prod_{r=1}^k|c_{i_rj_r}|^{1/2}
       \int_{T_z}(m(z+Vt)+\epsilon)^{-k/2}\,dt.
\end{equation}
\end{proposition}
\begin{proof}
The preimage of a test box is a rectangle
$\prod_r[\alpha_r,\beta_r]$, because the pairs are disjoint.
On the chord in direction $v_r$, the selected term has concave
quadratic coefficient $-\kappa_r$, where $\kappa_r=|c_{i_rj_r}|$.
Its envelope is no greater than the endpoint chord, so its gap at
$t_r$ is at least
$\kappa_r(t_r-\alpha_r)(\beta_r-t_r)$. Nonnegative termwise gaps
can be summed, giving on the owned feasible slice
\[
 m+\epsilon\geq\sum_r\kappa_r
                         (t_r-\alpha_r)(\beta_r-t_r).
\]
Arithmetic--geometric mean and the one-dimensional arcsine integral
bound each test's contribution to the integral by
$k^{-k/2}\pi^k\prod_r\kappa_r^{-1/2}$. Integrate over its owned subset
and bound by the full preimage rectangle. Degenerate rectangles have
zero $k$-measure. Summation proves the claim.
\end{proof}

For one such feasible chord through a minimizer,
$z_*+tv$ with $|t|\leq t_0$ and $t_0>0$, an upper growth bound
$m(z_*+tv)\leq Mt^2$ with $M>0$ yields a logarithmic lower bound by
\cref{geom:eq-matching}: the integral of $(Mt^2+\epsilon)^{-1/2}$ is
$2M^{-1/2}\operatorname{arsinh}(t_0\sqrt{M/\epsilon})$.
An unconstrained interior $C^{1,1}$ minimizer supplies this upper growth.
With constraints the feasible chord must be established separately.

\begin{proposition}[One-sided concave ray]\label{geom:prop-ray}
Suppose $z_*+tv$ is feasible for $0\leq t\leq t_0$, with $t_0>0$, a selected edge
satisfies $c_{ij}v_iv_j<0$, and
$\liminf_{t\downarrow0}m(z_*+tv)/t=0$. Then
$N_{\mathrm{cover}}(\epsilon)\to\infty$ as $\epsilon\downarrow0$,
including for constrained problems.
\end{proposition}
\begin{proof}
If covers with at most $N$ boxes existed along $\epsilon_\ell\downarrow0$,
their intersections with the ray would give at most $N$ closed
intervals covering $[0,t_0]$. Pad empty intersections by singleton
intervals and pass to a subsequence along which all endpoints converge.
The limit intervals still cover $[0,t_0]$, since the distance of any
point to their union is the limit of its distances to the finite unions.
For a limiting interval $[a,b]$ and $a<t<b$, endpoint interpolation
of the concave selected edge gives
\[
 m(z_*+tv)+\epsilon_\ell\geq
           |c_{ij}v_iv_j|(t-a_\ell)(b_\ell-t).
\]
For large $\ell$, $t$ lies in that interval. Passing to the limit
gives the same inequality at tolerance zero. Among finitely many
limit intervals covering points arbitrarily close to zero, one has
left endpoint zero and right endpoint $b>0$. Consequently
$m(z_*+tv)/t\geq|c_{ij}v_iv_j|(b-t)$ for $0<t<b$, contradicting the
liminf assumption.
\end{proof}

\begin{proposition}[A full aligned segment in two dimensions]
\label{geom:prop-aligned-two}
Let $f(x,y)=g(x,y)+cxy$, with $c\ne0$ and $g$ convex and continuous
on $[L_x,U_x]\times[L_y,U_y]$. If the entire segment
$\{a\}\times[L_y,U_y]$ is optimal, cutting at $x=a$ gives an exact
certificate with at most two leaves. If $a$ is an endpoint, the root
is exact. The assertion uses the termwise model
\cref{geom:eq-mcc-model} with no other constraints.
\end{proposition}

\Cref{app:geometry:proof-aligned-two} proves this statement using
convexity to identify the transverse slope along the segment. Its
full-length hypothesis matters. Alignment alone is insufficient in
higher dimensions, as the following example shows.

\begin{example}[An aligned three-dimensional segment]
\label{geom:ex-aligned-three}
On $D=F=[0,1]^3$, fix $a\in(0,1)$ and set
\[
 f=(x-a)^2+(x-a)(y-z)+(y-z)^2.
\]
Use the specified decomposition
$g=(x-a)^2+(y-z)^2-a(y-z)$ and $\phi=xy-xz$, with graph
$G=\{xy,xz\}$. The function $g$ is convex. Writing $X=x-a$ and
$Y=y-z$ gives
$f=X^2+XY+Y^2\geq(X^2+Y^2)/2$, so the optimum is zero and its
set is the segment $\{(a,t,t):0\leq t\leq1\}$.
Its tangent lies in the independent coordinates $\{y,z\}$ and its
line constant $\tau$ is zero. Nevertheless
\[
 N_{\mathrm{cover}}\asymp N_{\mathrm{rect}}
             \asymp N_{\mathrm{tree}}\asymp\epsilon^{-1/2}.
\]
For the lower bound consider the strip
$\{(x,t,t):|x-a|\leq h,\ 0\leq t\leq1\}$, with parameter area $2h$.
On it $m\leq h^2$. A valid test meets the strip in a parameter
rectangle. At its center the two edge lower bounds give gap at least
$w_xw_t/2$, so its area is at most $2(h^2+\epsilon)$.
Summing and taking $h=\sqrt\epsilon\leq\min(a,1-a)$ gives
$N_{\mathrm{cover}}\geq1/(2\sqrt\epsilon)$.

For the upper bound, a dyadic cube of side $s$ has
$\Gamma_B\leq s^2/2$. An unpruned cube therefore contains a point
with $|x-a|<s$ and $|y-z|<s$. At level $j$ at most four $x$-intervals
and at most $5\cdot2^j$ pairs of $y,z$ intervals can do so.
There are $O(2^j)$ unpruned cubes. All cubes are valid when
$s_j^2/2\leq\epsilon$. The geometric sum of active-parent counts
up to this level gives $O(\epsilon^{-1/2})$ processed nodes, and
coordinate cuts realize the dyadic construction. The fixed
decomposition thus gives diverging certificates even for this convex
quadratic objective.
\end{example}
